\section{Paradigmas modernos en la representación del espacio de búsqueda}

\subsection{De la nube de puntos a la variedad Topológica ($T^n$)}
Tradicionalmente, el Espacio de Planificación de Rutas (PPS) se definía analizando 
nubes de puntos o mallas discretas. Sin embargo, este enfoque sufre de 
\enquote{ceguera estructural} porque ignora las propiedades intrínsecas 
del espacio, y genera un costo computacional masivo por las constantes 
proyecciones requeridas para evitar obstáculos, exacerbado por la maldición 
de la dimensionalidad \parencite{radhakrishnan2024state, semwal2024manipulator}.

Para superar estas limitaciones, el paradigma actual modela el entorno como 
variedades topológicas (\textit{manifolds}), garantizando continuidad geométrica. 
Al tratar a los manipuladores robóticos como topológicamente homeomorfos a un 
toroide de $k$-dimensiones ($T^k$) estructurado por sus articulaciones 
\parencite{pierce2025constrained}, los planificadores pueden incrustar las 
rutas matemáticas usando cartas locales. Como resultado, la planificación 
navega de manera fluida y analítica a través de \enquote{pasajes estrechos} 
sin tener que reevaluar colisiones vértice por vértice \parencite{radhakrishnan2024state}.

\subsection{Mapeo de obstáculos mediante geometrías convexas}
Para optimizar aún más la detección de colisiones, la investigación 
migró de mallas poligonales a representaciones convexas. Enfoques como el 
\enquote{Mundo de Esferas} (\textit{Sphere World}) encapsulan los obstáculos 
dentro de uniones disjuntas de bolas cerradas, elipsoides o variedades carentes 
de bordes afilados (\textit{edge-free}) \parencite{radhakrishnan2024state}. El 
gran impacto de estas \enquote{geometrías de delimitación} es la reducción drástica 
de tiempo computacional, ya que permite que los algoritmos calculen analíticamente 
la divergencia más corta (la geodésica matemática) fluyendo alrededor de las 
fronteras de los objetos, garantizando seguridad y suavidad en la evasión.



\section{Análisis comparativo de algoritmos en espacios de configuración ($C$)}

\subsection{El puente algorítmico: Homotopía para evadir colisiones dinámicas}
Una vez que el espacio de búsqueda se ha mapeado analíticamente como una 
variedad ($T^n$), el siguiente desafío es enseñar al sistema qué rutas son 
matemáticamente equivalentes para saltar entre ellas si un obstáculo cierra el 
paso repentinamente. En este punto entran los métodos basados en homotopía 
(HBM) \parencite{radhakrishnan2024state2}.

La homotopía permite estructurar este amplio espacio agrupando infinitas 
configuraciones en \enquote{clases de homotopía}: familias de rutas que 
conectan dos puntos y que pueden deformarse elásticamente sin intersectarse con 
ningún obstáculo. Para clasificarlas y guiar al robot, enfoques modernos utilizan 
Firmas Matemáticas (\textit{H-signature}) o secuencias de cruces topológicos 
llamadas \enquote{Palabras canónicas} \parencite{Bhattacharya2012Topological}. 
La principal función de formar esta estructura es otorgar inmunidad a los 
planificadores ante los \enquote{mínimos locales}. Si un objeto dinámico bloquea 
un rumbo, un algoritmo homotópico avanzado como POSH rechaza la fricción de forzar 
un cruce y decide de forma inteligente \enquote{saltar} (\textit{switch}) 
intencionalmente hacia una familia de homotopía distinta (ej. rodear un bloque 
por la izquierda en lugar de la derecha).

Justamente, para poder construir estas ramas homotópicas y buscar la solución 
más óptima, la literatura confronta constantemente dos arquitecturas algorítmicas 
de base: los grafos de búsqueda matemática y el muestreo de exploración probabilístico.

\subsection{Eficiencia computacional: A* frente a RRT en variedades}
El desempeño entre planificadores basados en exploración por grafos y en 
muestreo aleatorio revela un compromiso crítico (\textit{trade-off}). En 
ambientes topológicos altamente desafiantes, los planificadores deterministas 
rigurosos como \textbf{HA*} aseguran la ruta matemáticamente óptima, pero a un 
costo temporal muy estricto, promediando 30.99 segundos y escalando exponencialmente 
al expandir árboles limitados por continuaciones de altas dimensiones 
\parencite{Bhattacharya2012Topological, radhakrishnan2024state}.

En marcado contraste, los enfoques probabilísticos abstractos como \textbf{HRRT} y 
variaciones híbridas como \textbf{AtlasRRT} son casi 130 veces más rápidos (con medias 
de 0.239 segundos en RRT). No obstante, esta tremenda agilidad compensa sacrificando 
fuertemente la calidad topológica original del trayecto, generando rutas de 1.4 a 
2.19 veces más largas e inestables que la propuesta matemática óptima 
\parencite{Bhattacharya2012Topological, radhakrishnan2024state}.

\subsection{Algoritmos de vanguardia en alta dimensionalidad}
Para escenarios con intensa estrechez estática se han propuesto métodos altamente 
resolutivos como \textbf{Many-RRT*}, una arquitectura que genera objetivos dinámicos 
paralelos para evadir singularidades computacionales cinemáticas. Al evaluarse 
rigurosamente en brazos manipuladores de 6 y 7 grados de libertad (UR10e y Panda) 
contra estándares base de muestreo (RRT*, RRT-Connect), este método logró una 
métrica disruptiva: una tasa de éxito casi perfecta del 100\% en laberintos 
caóticos donde todos los sistemas probabilísticos tradicionales colapsaron 
logrando tan sólo del 1.4\% al 5.0\% de éxito. Además de la agilidad de los 
caminos, logra estabilizar trayectorias mucho más uniformes convergiendo en 
soluciones útiles con resoluciones de convergencia 6 veces menores, permitiéndole 
reducir finalmente el costo métrico analítico global (distancia operada combinada 
en la variedad) hasta un 44.5\% gastando exactamente el mismo runtime operacional 
preasignado \parencite{belmont2026many}.

\section{Democratización del hardware y gemelos digitales}

\subsection{Democratización mediante bajo costo y su vacío pedagógico}
En los últimos años, la dependencia de la academia por los masivos y costosos 
robots industriales se mitigó por un avance sustancial con software Open-Source 
complementados por la fabricación aditiva (FDM 3D). Diseños como el manipulador 
articulado \textbf{G-ARM} han probado las bases pragmáticas en su despliegue 
educativo robótico integral mediante un costo estimado total de adquisición para 
los constructores de apenas 171 euros. Paralelamente, sistemas mecanum orientados a 
cinemáticas omnidireccionales diferenciales como \textbf{SMARTmBOT} limitan los gastos 
completos, proveyendo bases Raspberry Pi robustas y hasta formaciones LIDAR mediante 
Time Of Flight y cámaras visuales por aproximadamente 210 dólares netos, probando ser 
viables para replicaciones a gran escala en centros sin presupuesto elevado 
\parencite{vega2025g, jo2022smartmbot}.

Sin embargo, a pesar de este gran acierto algorítmico global, la gran crítica persiste 
ante estas mecánicas robóticas por fomentar el \enquote{vacío pedagógico} 
educativo \parencite{de2025development}. Las investigaciones de vanguardia 
delegan puramente las funciones complejas preenlatadas de estos manipuladores 
accesibles a \enquote{Cajas Negras} matemáticas de la industria (MoveIt2, OPML, 
algoritmos orientados como KDL) provocando que las mecánicas del sistema operen 
fluidamente en evasiones avanzadas, pero el diseñador/estudiante pierda tangibilidad. 
Estos métodos unificados blindan el conocimiento algebraico interno de manera que 
jamás se enfrentan transparentemente al modelado o la parametrización analítica de 
la propia topología en el $C-space$ limitando gravemente el entendimiento conceptual 
interno detrás de las planificaciones robóticas  
\parencite{semwal2024manipulator, vega2025g}.

\subsection{Cerrando el reality gap y asegurando la reproducibilidad}
Para cruzar la alta barrera aislacionista teórica matemática, métodos de 
instrumentación reciente interlazan de forma constante el mundo virtual mediante 
vinculaciones fluidas entre los planificadores y el hardware. Agrupados dentro de 
paradigmas como Gemelos Digitales asíncronos, se busca garantizar que las órbitas y 
resoluciones matemáticas colisiónen perfectamente entre la iteración teórica sin 
desviaciones abrumadoras causadas ante deficiencias motrices (resolviendo drásticamente 
el Reality-Gap de la fabricación) \parencite{de2025development}. 

No obstante, esta interactividad del software a menudo envejece críticamente mal. Las 
debilidades originadas por las dependencias binarias atadas a diferentes SO operativos 
Linux destrozan de raíz la reusabilidad del entorno cuando este expira, sumiendo los 
códigos a un ciclo corto. La respuesta a esta \textit{crisis de reproducibilidad} ha 
sido respaldada por \textbf{Docker}, permitiendo desplazar arquitecturas operativas 
fijadas permanentemente. En especial, mediante el uso de orquestadores como 
\textbf{Docker Compose} y el puenteo directo del servidor gráfico (X11), se logra 
aislar la renderización GUI y evitar la contaminación del sistema operativo local 
con librerías de un solo uso, operando entornos complejos mediante un simple comando 
de terminal unificada \parencite{martinez2022docker, cervera2023run}.

\section{Síntesis y brecha de investigación}

La revisión literaria evidencia una desconexión crítica actual. Por un lado, 
la matemática de las topologías teóricas ($T^k$) han alcanzado planificaciones 
inmejorables limitadas a centros con alto presupuesto o teorías exclusivas 
virtuales \parencite{radhakrishnan2024state}. Contrariamente, en la contraparte 
educativa del brazo robótico abierto, las aplicaciones son lideradas integralmente 
recargando los algoritmos a softwares caja-negra comerciales unificados 
\parencite{semwal2024manipulator, vega2025g}. La literatura explorada sugiere 
una falta de integración para documentar y construir desde su base algorítmica 
una propuesta de simulador educativo tipo \enquote{Caja Blanca}, exponiendo los 
cálculos topológicos de forma transparente y encapsulada para hardware 
comunitario \parencite{cervera2023run, de2025development}.

Basados en esta brecha del panorama académico referenciado, el argumento que 
fundamentaría y fortalecería esta propuesta se apoya de cuatro validaciones esenciales:

\begin{enumerate}
    \item \textbf{Transparencia frente a las \enquote{cajas negras}:} A diferencia de 
    las metodologías que ocultan la lógica en bloques automatizados como MoveIt2 
    \parencite{semwal2024manipulator}, operar con algoritmos A* o RRT construidos 
    desde cero ayuda a cerrar el \enquote{vacío pedagógico}, permitiendo al estudiante 
    interactuar directamente con la matemática del espacio de configuración topológico.
    \item \textbf{Mitigación al \textit{curve-jumping}:} A nivel metodológico 
    algebraico, el anclar el planificador A*/RRT resolviendo sobre métricas equivalentes 
    en $T^n$ blinda al algoritmo y lo hace inmune del crítico fenómeno en 
    \enquote{salto de curvas no-lineales} garantizando congruencias 
    \parencite{radhakrishnan2024state2, velez2022novel}.
    \item \textbf{Consistencia del entorno mediante Docker:} El uso de contenedores 
    elimina la acumulación de dependencias obsoletas y evita la modificación o 
    contaminación del sistema operativo local con software redundante. Esto asegura 
    que el código fuente actúe de manera inmune y replicable de manera inmediata sin 
    temor a la expiración de las versiones \parencite{cervera2023run}.
    \item \textbf{Validación física y reducción del reality gap:} Finalmente,
     aterrizar este enfoque contenedorizado sobre el diseño del brazo impreso 
     3D asíncrono permite validar que la teoría topológica planificada se traduzca 
     de forma fidedigna en los motores reales orientando una vinculación directa entre 
     la abstracción matemática y la práctica en la ingeniería mecatrónica 
     \parencite{de2025development}.
\end{enumerate}

En conclusión, este enfoque cruza y vincula tres áreas aisladas en la literatura: 
el modelado de la matemática pura de variedades topológicas, la difusión práctica 
pedagógica en la educación impulsada a coste libre, y las garantías de software 
modernas basadas en contenedores computacionales.